% GENERATED FILE. Edit canonical lessons, then run npm run generate:books.
% canonical-source-hash: fnv1a64:4a68ae484e03f3ca
% canonical-lessons: BN-C469-ghama, BN-C469-bron, BN-C469-til, BN-C469-dag, BN-C469-lom, BN-R469-review-462-469

\chapter{Sweat and Skin}
\label{ch:bn-sweat-and-skin}
\hlchaptermodality{469}

\begin{chapteropening}
\textbf{By the end of this chapter:} \emph{I can talk about sweat, spots, moles and stains.}

Complete the last lesson of chapter 469: I can talk about sweat, spots, moles and stains.
\end{chapteropening}

\section[\texorpdfstring{\bn{ঘামা} (ghāmā)}{ঘামা (ghāmā)}]{\bn{ঘামা} (ghāmā) — to sweat}

\label{lesson:BN-C469-ghama}



\begin{quote}
Before the new one: say the Bengali for a vein, then the Bengali for to thump.
\end{quote}

\subsection*{You'll want to know: \bn{ঘামা}}
\textbf{\bn{ঘামা}} — \emph{ghāmā} — \textquotedblleft{}to sweat\textquotedblright{}.

\begin{quote}
\bn{গরমে} \bn{বাসে} \bn{বসে} \bn{আমি} \bn{ঘামছিলাম।}
\end{quote}

\begin{quote}
\emph{gôrôme bāse bôse āmi ghāmchhilām.}
\end{quote}

\begin{quote}
Sitting on the bus in the heat, I was sweating.
\end{quote}

Past (he/she …): \bn{তিনি} \bn{রোদে} \bn{ঘামলেন}, \emph{tini rode ghāmlen}.

The noun is \bn{ঘাম}, sweat: \bn{ঘামে} \bn{ভেজা} \bn{জামা}, a shirt wet with sweat.

\begin{cousinweb}
\textbf{\bn{ঘামা}} (\emph{ghāmā}) is the verb of \emph{ghām}, \textquotedblleft{}sweat\textquotedblright{}, which CDIAL lists for Bengali under Sanskrit \emph{gharma}, \textquotedblleft{}heat\textquotedblright{} (CDIAL 4445).
\end{cousinweb}

\begin{grammarlens}[title={What you've built}]
One more word for describing what you lived and felt.
\end{grammarlens}

\subsection*{Your turn}
\begin{itemize}
\raggedright
  \item \emph{Say it:} \emph{ghāmā}
  \item \emph{Say it:} \emph{ghāmā}, once more
  \item \emph{Say it:} \emph{gôrôme bāse bôse āmi ghāmchhilām.} — Sitting on the bus in the heat, I was sweating.
  \item \emph{Recall:} say the Bengali for a vein, then the Bengali for to thump, then say \emph{ghāmā} again
\end{itemize}

\begin{tcolorbox}[breakable,colback=teal!4,colframe=teal!35!black,title={Before you move on}]
What does \bn{ঘামা} mean? (\textquotedblleft{}To sweat\textquotedblright{}.) Say it once more.
\end{tcolorbox}
\section[\texorpdfstring{\bn{ব্রণ} (brôn)}{ব্রণ (brôn)}]{\bn{ব্রণ} (brôn) — a pimple; acne}

\label{lesson:BN-C469-bron}



\begin{quote}
Before the new one: say the Bengali for to thump, then the Bengali for to sweat.
\end{quote}

\subsection*{You'll want to know: \bn{ব্রণ}}
\textbf{\bn{ব্রণ}} — \emph{brôn} — \textquotedblleft{}a pimple; acne\textquotedblright{}.

\begin{quote}
\bn{মুখে} \bn{ব্রণ} \bn{হলে} \bn{অরুণ} \bn{লজ্জা} \bn{পায়।}
\end{quote}

\begin{quote}
\emph{mukhe brôn hôle ôrun lôjjā pāy.}
\end{quote}

\begin{quote}
When he gets spots on his face, Arun feels shy.
\end{quote}

\begin{cousinweb}
\textbf{\bn{ব্রণ}} (\emph{brôn}) is a tatsama, Sanskrit \emph{vraṇa}, \textquotedblleft{}an abscess, a boil, a scar\textquotedblright{} (CDIAL 12227); Bengali has narrowed it to a pimple.
\end{cousinweb}

\begin{grammarlens}[title={What you've built}]
One more word for describing what you lived and felt.
\end{grammarlens}

\subsection*{Your turn}
\begin{itemize}
\raggedright
  \item \emph{Say it:} \emph{brôn}
  \item \emph{Say it:} \emph{brôn}, once more
  \item \emph{Say it:} \emph{mukhe brôn hôle ôrun lôjjā pāy.} — When he gets spots on his face, Arun feels shy.
  \item \emph{Recall:} say the Bengali for to thump, then the Bengali for to sweat, then say \emph{brôn} again
\end{itemize}

\begin{tcolorbox}[breakable,colback=teal!4,colframe=teal!35!black,title={Before you move on}]
What does \bn{ব্রণ} mean? (\textquotedblleft{}A pimple; acne\textquotedblright{}.) Say it once more.
\end{tcolorbox}
\section[\texorpdfstring{\bn{তিল} (til)}{তিল (til)}]{\bn{তিল} (til) — a mole (on the skin)}

\label{lesson:BN-C469-til}



\begin{quote}
Before the new one: say the Bengali for to sweat, then the Bengali for a pimple.
\end{quote}

\subsection*{You'll want to know: \bn{তিল}}
\textbf{\bn{তিল}} — \emph{til} — \textquotedblleft{}a mole (on the skin)\textquotedblright{}.

\begin{quote}
\bn{মীরার} \bn{গালে} \bn{একটা} \bn{ছোট} \bn{তিল} \bn{আছে।}
\end{quote}

\begin{quote}
\emph{mīrār gāle ekṭā chhoṭo til āchhe.}
\end{quote}

\begin{quote}
Mira has a small mole on her cheek.
\end{quote}

\bn{তিল} is also the sesame seed; CDIAL thinks the two may be one word.

\begin{cousinweb}
\textbf{\bn{তিল}} (\emph{til}) is Sanskrit \emph{tila}, \textquotedblleft{}a mole on the skin\textquotedblright{} (CDIAL 5828), which CDIAL thinks may be the same word as \emph{tila}, the sesame seed: a seed-sized dark spot.
\end{cousinweb}

\begin{grammarlens}[title={What you've built}]
One more word for describing what you lived and felt.
\end{grammarlens}

\subsection*{Your turn}
\begin{itemize}
\raggedright
  \item \emph{Say it:} \emph{til}
  \item \emph{Say it:} \emph{til}, once more
  \item \emph{Say it:} \emph{mīrār gāle ekṭā chhoṭo til āchhe.} — Mira has a small mole on her cheek.
  \item \emph{Recall:} say the Bengali for to sweat, then the Bengali for a pimple, then say \emph{til} again
\end{itemize}

\begin{tcolorbox}[breakable,colback=teal!4,colframe=teal!35!black,title={Before you move on}]
What does \bn{তিল} mean? (\textquotedblleft{}A mole (on the skin)\textquotedblright{}.) Say it once more.
\end{tcolorbox}
\section[\texorpdfstring{\bn{দাগ} (dāg)}{দাগ (dāg)}]{\bn{দাগ} (dāg) — a mark; a stain; a scar}

\label{lesson:BN-C469-dag}



\begin{quote}
Before the new one: say the Bengali for a pimple, then the Bengali for a mole (on the skin).
\end{quote}

\subsection*{You'll want to know: \bn{দাগ}}
\textbf{\bn{দাগ}} — \emph{dāg} — \textquotedblleft{}a mark; a stain; a scar\textquotedblright{}.

\begin{quote}
\bn{জামায়} \bn{চায়ের} \bn{দাগ} \bn{লাগল।}
\end{quote}

\begin{quote}
\emph{jāmāy chāyer dāg lāglô.}
\end{quote}

\begin{quote}
Tea left a stain on the shirt.
\end{quote}

\bn{কাটা} \bn{দাগ}, the scar of a cut.

\begin{cousinweb}
\textbf{\bn{দাগ}} (\emph{dāg}) is usually traced to Persian \emph{dāgh}, \textquotedblleft{}a brand, a mark\textquotedblright{}.
\end{cousinweb}

\begin{grammarlens}[title={What you've built}]
One more word for describing what you lived and felt.
\end{grammarlens}

\subsection*{Your turn}
\begin{itemize}
\raggedright
  \item \emph{Say it:} \emph{dāg}
  \item \emph{Say it:} \emph{dāg}, once more
  \item \emph{Say it:} \emph{jāmāy chāyer dāg lāglô.} — Tea left a stain on the shirt.
  \item \emph{Recall:} say the Bengali for a pimple, then the Bengali for a mole (on the skin), then say \emph{dāg} again
\end{itemize}

\begin{tcolorbox}[breakable,colback=teal!4,colframe=teal!35!black,title={Before you move on}]
What does \bn{দাগ} mean? (\textquotedblleft{}A mark; a stain; a scar\textquotedblright{}.) Say it once more.
\end{tcolorbox}
\section[\texorpdfstring{\bn{লোম} (lom)}{লোম (lom)}]{\bn{লোম} (lom) — body hair; fur}

\label{lesson:BN-C469-lom}



\begin{quote}
Before the new one: say the Bengali for a mole (on the skin), then the Bengali for a stain.
\end{quote}

\subsection*{You'll want to know: \bn{লোম}}
\textbf{\bn{লোম}} — \emph{lom} — \textquotedblleft{}body hair; fur\textquotedblright{}.

\begin{quote}
\bn{ভয়ে} \bn{আমার} \bn{হাতের} \bn{লোম} \bn{খাড়া} \bn{হয়ে} \bn{গেল।}
\end{quote}

\begin{quote}
\emph{bhôye āmār hāter lom khāṛā hôye gelô.}
\end{quote}

\begin{quote}
The hairs on my arms stood on end with fear.
\end{quote}

\begin{cousinweb}
\textbf{\bn{লোম}} (\emph{lom}) is a tatsama, Sanskrit \emph{loman}, \textquotedblleft{}hair on the body of animals or people\textquotedblright{} (CDIAL 11154).
\end{cousinweb}

\begin{grammarlens}[title={What you've built}]
One more word for describing what you lived and felt.
\end{grammarlens}

\subsection*{Your turn}
\begin{itemize}
\raggedright
  \item \emph{Say it:} \emph{lom}
  \item \emph{Say it:} \emph{lom}, once more
  \item \emph{Say it:} \emph{bhôye āmār hāter lom khāṛā hôye gelô.} — The hairs on my arms stood on end with fear.
  \item \emph{Recall:} say the Bengali for a mole (on the skin), then the Bengali for a stain, then say \emph{lom} again
\end{itemize}

\begin{tcolorbox}[breakable,colback=teal!4,colframe=teal!35!black,title={Before you move on}]
What does \bn{লোম} mean? (\textquotedblleft{}Body hair; fur\textquotedblright{}.) Say it once more.
\end{tcolorbox}
\section[(dialogue)]{Looking back: from Lazy and Industrious to Sweat and Skin}

\label{lesson:BN-R469-review-462-469}



\begin{quote}
Say the Bengali for lazy and industrious.
\end{quote}

\subsection*{Your turn}
Say these aloud:

\begin{itemize}
\raggedright
  \item lazy, industrious, in disorder, to shirk, to dawdle
  \item naughty, witty, talkative, to tease, to tickle
  \item grave, unpretentious, fussy, indifferent, to adapt
  \item conscience, judicious, gifted, grateful, to weigh up
  \item the body, an arm, a finger, the back, the nape
  \item a lip, the jaw, the gums, the voice, to knit one's brows
  \item the heart (the organ), the liver, the stomach (the organ), a vein, to thump
  \item to sweat, a pimple, a mole (on the skin), a stain, body hair
\end{itemize}

\begin{tcolorbox}[breakable,colback=teal!4,colframe=teal!35!black,title={Before you move on}]
Lazy? (\textbf{\bn{অলস}}, \emph{ôlôs}.) Industrious? (\textbf{\bn{কর্মঠ}}, \emph{kôrmôṭh}.) In disorder? (\textbf{\bn{এলোমেলো}}, \emph{elomelo}.) To shirk? (\textbf{\bn{ফাঁকি} \bn{দেওয়া}}, \emph{phā̃ki deoyā}.) To dawdle? (\textbf{\bn{গড়িমসি} \bn{করা}}, \emph{gôṛimôsi kôrā}.) Naughty? (\textbf{\bn{দুষ্টু}}, \emph{dushṭu}.) Witty? (\textbf{\bn{রসিক}}, \emph{rôsik}.) Talkative? (\textbf{\bn{বাচাল}}, \emph{bāchāl}.) To tease? (\textbf{\bn{ঠাট্টা} \bn{করা}}, \emph{ṭhāṭṭā kôrā}.) To tickle? (\textbf{\bn{কাতুকুতু} \bn{দেওয়া}}, \emph{kātukutu deoyā}.) Grave? (\textbf{\bn{গম্ভীর}}, \emph{gômbhīr}.) Unpretentious? (\textbf{\bn{সাদাসিধে}}, \emph{sādāsidhe}.) Fussy? (\textbf{\bn{খুঁতখুঁতে}}, \emph{khũtkhũte}.) Indifferent? (\textbf{\bn{উদাসীন}}, \emph{udāsīn}.) To adapt? (\textbf{\bn{মানিয়ে} \bn{নেওয়া}}, \emph{māniye neoyā}.) Conscience? (\textbf{\bn{বিবেক}}, \emph{bibek}.) Judicious? (\textbf{\bn{বিচক্ষণ}}, \emph{bichôkkhôn}.) Gifted? (\textbf{\bn{মেধাবী}}, \emph{medhābī}.) Grateful? (\textbf{\bn{কৃতজ্ঞ}}, \emph{kritôgyô}.) To weigh up? (\textbf{\bn{বিবেচনা} \bn{করা}}, \emph{bibechônā kôrā}.) The body? (\textbf{\bn{শরীর}}, \emph{shôrīr}.) An arm? (\textbf{\bn{বাহু}}, \emph{bāhu}.) A finger? (\textbf{\bn{আঙুল}}, \emph{āngul}.) The back? (\textbf{\bn{পিঠ}}, \emph{piṭh}.) The nape? (\textbf{\bn{ঘাড়}}, \emph{ghāṛ}.) A lip? (\textbf{\bn{ঠোঁট}}, \emph{ṭhõṭ}.) The jaw? (\textbf{\bn{চোয়াল}}, \emph{choyāl}.) The gums? (\textbf{\bn{মাড়ি}}, \emph{māṛi}.) The voice? (\textbf{\bn{কণ্ঠ}}, \emph{kônṭhô}.) To knit one's brows? (\textbf{\bn{কোঁচকানো}}, \emph{kõchkāno}.) The heart (the organ)? (\textbf{\bn{হৃৎপিণ্ড}}, \emph{hritpinḍô}.) The liver? (\textbf{\bn{যকৃৎ}}, \emph{jôkrit}.) The stomach (the organ)? (\textbf{\bn{পাকস্থলী}}, \emph{pākôsthôlī}.) A vein? (\textbf{\bn{শিরা}}, \emph{shirā}.) To thump? (\textbf{\bn{ধুকপুক} \bn{করা}}, \emph{dhukpuk kôrā}.) To sweat? (\textbf{\bn{ঘামা}}, \emph{ghāmā}.) A pimple? (\textbf{\bn{ব্রণ}}, \emph{brôn}.) A mole (on the skin)? (\textbf{\bn{তিল}}, \emph{til}.) A stain? (\textbf{\bn{দাগ}}, \emph{dāg}.) Body hair? (\textbf{\bn{লোম}}, \emph{lom}.)
\end{tcolorbox}
